\lecture{5}{2026-09-24}{Primality Testing}

Public-key cryptosystems, such as the Diffie--Hellman key exchange and RSA, rely on the generation of large prime numbers.
Because factoring large integers is computationally intractable while multiplying them is efficient, finding large primes is a fundamental primitive.
The standard strategy is to sample random integers of the desired bit-length and test them for primality.

\section*{Prime Density and Prime Generation}

To generate primes by random sampling, primes must be sufficiently dense.
Let \(\pi(n)\) denote the number of primes in the interval \([1, n]\).
The asymptotic distribution of primes is governed by the \emph{Prime Number Theorem}.

\begin{theorem}[Prime Number Theorem]
    As \(n \to \infty\),
    \[
        \lim_{n \to \infty} \frac{\pi(n)}{n / \ln n} = 1.
    \]
\end{theorem}

Explicit bounds are known for finite intervals. For example, Rosser~\cite{Rosser1941} proved that for all \(n \geq 55\),
\[
    \frac{n}{\ln n + 2} < \pi(n) < \frac{n}{\ln n - 4}.
\]
Using these bounds, a random integer of bit-length \(\ell\) (in the interval \([2^{\ell-1}, 2^\ell]\)) is prime with probability at least \(\frac{1}{\ell \ln 2 + 2}\).
While the expected number of trials to find a prime is \(O(\ell)\), sampling \(\ell^2\) candidate integers independently at random ensures that the probability of failing to find any prime is at most
\[
    \left( 1 - \frac{1}{\ell \ln 2 + 2} \right)^{\ell^2} \leq e^{-\ell^2 / (\ell \ln 2 + 2)} \leq e^{-\ell / (\ln 2 + 2)}.
\]
which is negligible in the bit-length \(\ell\).
Hence, to generate large primes, it suffices to have an efficient algorithm to test whether a given candidate integer is prime.

While a deterministic polynomial-time algorithm exists (the AKS test by Agrawal, Kayal, and Saxena~\cite{AgrawalKayalSaxena2004}), its high complexity makes it impractical for cryptographically large primes.
In practice, fast probabilistic tests such as the Miller--Rabin test are used instead.

\section*{The Fermat Primality Test}

A natural starting point is Fermat's Little Theorem.

\begin{theorem}[Fermat's Little Theorem]
    If \(p\) is prime, then for every \(a \in \mathbb{Z}\) with \(\gcd(a, p) = 1\),
    \[
        a^{p-1} \equiv 1 \pmod{p}.
    \]
\end{theorem}

If an integer \(n\) is composite, we might hope that \(a^{n-1} \not\equiv 1 \pmod{n}\) for most choices of \(a \in \mathbb{Z}_n^*\).
This leads to the \emph{Fermat Primality Test}.

\begin{algorithm}[htbp]
    \caption{Fermat Primality Test}\label{alg:fermat-test}
    \begin{algorithmic}[1]
        \Input integer \(n \geq 3\)
        \State choose a random nonzero \(a \in \mathbb{Z}_n\)
        \If{\(\gcd(a, n) > 1\)}
            \State \Return \textup{composite}
        \EndIf
        \If{\(a^{n-1} \not\equiv 1 \pmod{n}\)}
            \State \Return \textup{composite}
        \Else
            \State \Return \textup{probably prime}
        \EndIf
    \end{algorithmic}
\end{algorithm}

For a random nonzero \(a \in \mathbb{Z}_n\), the condition \(\gcd(a, n) > 1\) occurs with negligible probability when \(n\) has large prime factors.
Hence, the efficacy of the test hinges on whether \(a^{n-1} \not\equiv 1 \pmod{n}\) occurs frequently among units in \(\mathbb{Z}_n^*\).

\begin{proposition}\label{prop:fermat-subgroup}
    Let \(n\) be a composite integer.
    If there exists at least one base \(a_0 \in \mathbb{Z}_n^*\) such that \(a_0^{n-1} \not\equiv 1 \pmod{n}\), then
    \[
        a^{n-1} \not\equiv 1 \pmod{n}
    \]
    for at least half of all elements \(a \in \mathbb{Z}_n^*\).
\end{proposition}

\begin{proof}
    Consider the set \( S = \{a \in \mathbb{Z}_n^* : a^{n-1} \equiv 1 \pmod{n}\}\).
    Since \(1 \in S\) and \(S\) is closed under multiplication and inversion in \(\mathbb{Z}_n^*\), \(S\) is a subgroup of the multiplicative group \(\mathbb{Z}_n^*\).
    By assumption, \(a_0 \in \mathbb{Z}_n^* \setminus S\), so \(S\) is a proper subgroup of \(\mathbb{Z}_n^*\).
    By Lagrange's theorem, \(|S|\) divides \(|\mathbb{Z}_n^*|\), which implies \( |S| \leq \frac{1}{2} |\mathbb{Z}_n^*|\).
    Therefore, at least half of the elements in \(\mathbb{Z}_n^*\) satisfy \(a^{n-1} \not\equiv 1 \pmod{n}\).
\end{proof}

\begin{definition}[Fermat Witness and Fermat Liar]
    Let \(n\) be an odd composite integer and \(a \in \mathbb{Z}_n^*\).
    \begin{itemize}
        \item If \(a^{n-1} \not\equiv 1 \pmod{n}\), then \(a\) is called a \emph{Fermat witness} for the compositeness of \(n\).
        \item If \(a^{n-1} \equiv 1 \pmod{n}\), then \(a\) is called a \emph{Fermat liar} for \(n\).
    \end{itemize}
\end{definition}

\Cref{prop:fermat-subgroup} implies that if a single Fermat witness exists in \(\mathbb{Z}_n^*\), a randomly selected base \(a \in \mathbb{Z}_n^*\) reveals compositeness with probability at least \(1/2\).
Unfortunately, there exist composite numbers with no Fermat witnesses at all in \(\mathbb{Z}_n^*\), known as \emph{Carmichael numbers}.

\begin{definition}[Carmichael Number]
    A composite integer \(n\) is called a \emph{Carmichael number} if every element \(a \in \mathbb{Z}_n^*\) is a Fermat liar, that is,
    \[
        a^{n-1} \equiv 1 \pmod{n} \quad \text{for all } a \in \mathbb{Z}_n^*.
    \]
\end{definition}

\begin{example}
    The smallest Carmichael number is \(561 = 3 \cdot 11 \cdot 17\).
    By Korselt's criterion, a composite integer \(n\) is a Carmichael number if and only if \(n\) is square-free and \((p - 1) \mid (n - 1)\) for every prime divisor \(p\) of \(n\).
    For \(n = 561\), we have \(n - 1 = 560\), which is divisible by:
    \[
        3 - 1 = 2, \quad 11 - 1 = 10, \quad \text{and} \quad 17 - 1 = 16.
    \]
    Thus, for every \(a \in \mathbb{Z}_{561}^*\), Fermat's Little Theorem gives \(a^{560} \equiv 1\) modulo \(3\), \(11\), and \(17\), so by the Chinese Remainder Theorem,
    \[
        a^{560} \equiv 1 \pmod{561}.
    \]
\end{example}

\begin{remark}
    It is known that for any \(n \in \mathbb{Z}\), the number of Carmichael numbers in the interval \([1, n]\) is at least \(n^{2/7}\) (see~\cite{Harman2005}; see also~\cite{AlfordGranvillePomerance1994}).
    For this reason, the Fermat test does not always reliably distinguish between prime numbers and composite numbers, since there exist infinitely many composite numbers for which the Fermat test almost always outputs ``probably prime''.
\end{remark}

\section*{The Miller--Rabin Primality Test}

The Miller--Rabin test strengthens the Fermat test by exploiting the fact that over a field, the equation \(x^2 = 1\) has only the trivial roots \(x = \pm 1\).

\begin{lemma}\label{lem:roots-of-unity}
    If \(p\) is prime and \(x \in \mathbb{Z}_p\) satisfies \(x^2 \equiv 1 \pmod{p}\), then \(x \equiv \pm 1 \pmod{p}\).
\end{lemma}

\begin{proof}
    The condition \(x^2 \equiv 1 \pmod{p}\) is equivalent to \(p \mid (x^2 - 1) = (x - 1)(x + 1)\).
    By Euclid's lemma, since \(p\) is prime, \(p \mid (x - 1)\) or \(p \mid (x + 1)\), which means \(x \equiv 1 \pmod{p}\) or \(x \equiv -1 \pmod{p}\).
\end{proof}

If \(n\) is composite and has at least two distinct odd prime factors, the Chinese Remainder Theorem implies that \(x^2 \equiv 1 \pmod{n}\) has at least \(2^2 = 4\) distinct solutions modulo \(n\), creating \emph{non-trivial square roots of \(1\)}.
The Miller--Rabin test hunts for such non-trivial roots in the sequence of powers leading up to \(a^{n-1}\).

Let \(n \geq 3\) be an odd integer. We factor \(n - 1\) as
\[
    n - 1 = 2^e k, \quad \text{where \(k\) is odd and \(e \geq 1\)}.
\]
For any \(a \in \mathbb{Z}_n^*\), consider the sequence obtained by repeated squaring:
\[
    a^k, \; a^{2k}, \; a^{2^2 k}, \; \dots, \; a^{2^{e-1} k}, \; a^{2^e k} = a^{n-1} \pmod{n}.
\]

\begin{theorem}[Miller's Condition~\cite{Miller1976}]\label{thm:miller-condition}
    Let \(n \geq 3\) be an odd prime and write \(n - 1 = 2^e k\) with \(k\) odd.
    For every \(a \in \mathbb{Z}_n^*\), at least one of the following holds:
    \begin{enumerate}[label=\textup{(\roman*)}]
        \item \(a^k \equiv 1 \pmod{n}\);
        \item there exists some \(r \in \{0, 1, \dots, e-1\}\) such that \(a^{2^r k} \equiv -1 \pmod{n}\).
    \end{enumerate}
\end{theorem}

\begin{proof}
    Since \(n\) is prime, Fermat's Little Theorem guarantees that the last term of the sequence is \(a^{2^e k} = a^{n-1} \equiv 1 \pmod{n}\).
    If \(a^k \equiv 1 \pmod{n}\), condition (i) holds.
    Otherwise, let \(r \in \{0, 1, \dots, e-1\}\) be the largest index such that \(a^{2^r k} \not\equiv 1 \pmod{n}\).
    Then the next term in the sequence is \(a^{2^{r+1} k} = (a^{2^r k})^2 \equiv 1 \pmod{n}\).
    Thus, \(x = a^{2^r k}\) is a square root of \(1\) modulo \(n\) that is not congruent to \(1\).
    By \Cref{lem:roots-of-unity}, the only square roots of \(1\) modulo a prime are \(\pm 1\), so we must have \(a^{2^r k} \equiv -1 \pmod{n}\), satisfying condition (ii).
\end{proof}

\begin{definition}[Strong Witness and Strong Liar]
    Let \(n \geq 3\) be an odd composite integer with \(n - 1 = 2^e k\) (\(k\) odd), and let \(a \in \mathbb{Z}_n^*\).
    \begin{itemize}
        \item We say \(a\) is a \emph{strong witness} for \(n\) if
        \[
            a^k \not\equiv 1 \pmod{n} \quad \text{and} \quad a^{2^r k} \not\equiv -1 \pmod{n} \quad \text{for all } r \in \{0, 1, \dots, e - 1\}.
        \]
        \item Otherwise, \(a\) is called a \emph{strong liar} for \(n\) (and \(n\) is said to be a \emph{strong pseudoprime} to base \(a\)).
    \end{itemize}
\end{definition}

If \(a\) is a strong witness, then either \(a^{n-1} \not\equiv 1 \pmod{n}\) (a Fermat witness) or we encounter some \(x \not\equiv \pm 1 \pmod{n}\) with \(x^2 \equiv 1 \pmod{n}\), which yields a non-trivial square root of \(1\) modulo \(n\).
In either case, \(n\) is definitively composite.

\begin{algorithm}[htbp]
    \caption{Miller--Rabin Primality Test}\label{alg:miller-rabin}
    \begin{algorithmic}[1]
        \Input odd integer \(n \geq 3\) and candidate base \(a \in \{2, \dots, n-2\}\)
        \Output \textup{composite} or \textup{probably prime}
        \State factor \(n - 1 = 2^e k\) with \(k\) odd
        \State set \(x \gets a^k \pmod{n}\)
        \If{\(x \equiv 1 \pmod{n}\) \textbf{or} \(x \equiv -1 \pmod{n}\)}
            \State \Return \textup{probably prime} \Comment{\(a\) is a strong liar}
        \EndIf
        \For{\(r = 1\) \textbf{to} \(e - 1\)}
            \State set \(x \gets x^2 \pmod{n}\)
            \If{\(x \equiv -1 \pmod{n}\)}
                \State \Return \textup{probably prime} \Comment{\(a\) is a strong liar}
            \EndIf
            \If{\(x \equiv 1 \pmod{n}\)}
                \State \Return \textup{composite} \Comment{found non-trivial square root of \(1\)}
            \EndIf
        \EndFor
        \State \Return \textup{composite} \Comment{\(a^{n-1} \not\equiv 1 \pmod{n}\)}
    \end{algorithmic}
\end{algorithm}

Unlike the Fermat test, there are no composite numbers that defeat the Miller--Rabin test for all bases: even Carmichael numbers possess many strong witnesses.

\begin{theorem}\label{thm:strong-liars-bound}
    Let \(n \geq 3\) be an odd composite integer.
    Then at most \(\frac{1}{2}\) of the elements in \(\mathbb{Z}_n^*\) are strong liars for \(n\).
\end{theorem}

\begin{proof}
    We divide the proof into two cases depending on whether \(n\) is a prime power.

    \paragraph{Case 1: \(n = p^m\) for some prime \(p\) and \(m \geq 2\).}
    If \(a \in \mathbb{Z}_n^*\) is a strong liar, then \(a^{n-1} \equiv 1 \pmod{n}\).
    By Euler's theorem, we also have \(a^{\phi(n)} \equiv 1 \pmod{n}\), where \(\phi(n) = p^{m-1}(p-1)\).
    The order of \(a\) in \(\mathbb{Z}_n^*\) must divide both \(n - 1 = p^m - 1\) and \(\phi(n) = p^{m-1}(p-1)\).
    Since \(\gcd(p^m - 1, p^{m-1}) = 1\),
    \[
        \gcd(p^m - 1, p^{m-1}(p-1)) = \gcd(p^m - 1, p - 1) = p - 1.
    \]
    Thus, every strong liar must satisfy \(a^{p-1} \equiv 1 \pmod{n}\).
    The set \(H = \{a \in \mathbb{Z}_n^* : a^{p-1} \equiv 1 \pmod{n}\}\) forms a subgroup of \(\mathbb{Z}_n^*\).
    To show that \(H\) is proper, consider \(a = 1 + p^{m-1}\).
    By the binomial theorem,
    \[
        (1 + p^{m-1})^{p-1} = 1 + (p - 1)p^{m-1} + \sum_{j=2}^{p-1} \binom{p-1}{j} p^{j(m-1)} \equiv 1 - p^{m-1} \not\equiv 1 \pmod{p^m},
    \]
    since \(m \geq 2\) and \(2(m-1) \geq m\).
    Thus \(1 + p^{m-1} \notin H\), so \(H\) is a proper subgroup of \(\mathbb{Z}_n^*\).
    By Lagrange's theorem, \(|H| \leq \frac{1}{2} |\mathbb{Z}_n^*|\), so at most half of the elements in \(\mathbb{Z}_n^*\) are strong liars.

    \paragraph{Case 2: \(n\) is not a prime power.}
    Since \(n\) is not a prime power and is odd, we can write \(n = p^m q\) where \(p\) is an odd prime, \(m \geq 1\), and \(\gcd(p^m, q) = 1\) with \(q > 1\).
    Let \(i_0\) be the maximal index in \(\{0, 1, \dots, e-1\}\) for which there exists an element \(a_0 \in \mathbb{Z}_n^*\) satisfying
    \[
        a_0^{2^{i_0} k} \equiv -1 \pmod{n}.
    \]
    Such an index exists because \((-1)^k = -1\) (as \(k\) is odd), so \(i = 0\) always satisfies this with \(a = -1\).
    Now define the set
    \[
        G = \{a \in \mathbb{Z}_n^* : a^{2^{i_0} k} \equiv \pm 1 \pmod{n}\}.
    \]
    The set \(G\) is a subgroup of \(\mathbb{Z}_n^*\).
    Moreover, every strong liar \(a\) lies in \(G\):
    \begin{itemize}
        \item If \(a^k \equiv 1 \pmod{n}\), then \(a^{2^{i_0} k} \equiv 1 \pmod{n}\), so \(a \in G\).
        \item If \(a^{2^r k} \equiv -1 \pmod{n}\) for some \(r \leq i_0\), then \(a^{2^{i_0} k} \equiv (a^{2^r k})^{2^{i_0 - r}} \equiv (-1)^{2^{i_0 - r}} \equiv \pm 1 \pmod{n}\), so \(a \in G\).
    \end{itemize}
    We now show that \(G\) is a proper subgroup of \(\mathbb{Z}_n^*\).
    By the Chinese Remainder Theorem, choose \(b \in \mathbb{Z}_n^*\) such that
    \[
        b \equiv a_0 \pmod{p^m} \quad \text{and} \quad b \equiv 1 \pmod{q}.
    \]
    Raising to the power \(2^{i_0} k\), we obtain
    \[
        b^{2^{i_0} k} \equiv a_0^{2^{i_0} k} \equiv -1 \pmod{p^m}, \quad \text{and} \quad b^{2^{i_0} k} \equiv 1^{2^{i_0} k} \equiv 1 \pmod{q}.
    \]
    Since \(p^m \geq 3\) and \(q \geq 3\), \(-1 \not\equiv 1\) modulo \(p^m\) and modulo \(q\).
    For \(b^{2^{i_0} k}\) to equal \(1 \pmod{n}\), it would have to be \(1 \pmod{p^m}\); for it to equal \(-1 \pmod{n}\), it would have to be \(-1 \pmod{q}\).
    Since neither holds, \(b^{2^{i_0} k} \not\equiv \pm 1 \pmod{n}\), and hence \(b \notin G\).
    Therefore \(G\) is a proper subgroup of \(\mathbb{Z}_n^*\), and by Lagrange's theorem,
    \[
        |G| \leq \frac{1}{2} |\mathbb{Z}_n^*|.
    \]
    Since all strong liars belong to \(G\), at most \(\frac{1}{2}|\mathbb{Z}_n^*|\) elements are strong liars.
\end{proof}

\begin{remark}[Rabin's Bound]
    A more refined analysis due to Rabin~\cite{Rabin1980} shows that for any odd composite \(n\), the number of strong liars is actually bounded by
    \[
        |\{\text{strong liars in } \mathbb{Z}_n^*\}| \leq \frac{1}{4} \phi(n).
    \]
\end{remark}

\begin{remark}[Error Probability and Complexity]
    If \(n\) is composite and we run the Miller--Rabin test with \(t\) independent, uniformly chosen bases \(a_1, \dots, a_t \in \{2, \dots, n-2\}\), the probability that \(n\) is declared ``probably prime'' in all \(t\) iterations is at most
    \[
        \left( \frac{1}{4} \right)^t = 2^{-2t}.
    \]
    Setting \(t = 64\) guarantees a false positive probability less than \(2^{-128}\), which is cryptographically negligible.
    
    Each round requires computing \(a^k \pmod{n}\) and at most \(e - 1\) modular squarings, for a total of \(O(\log n)\) modular multiplications.
    With standard \(O(\log^2 n)\) arithmetic, each round runs in \(O(\log^3 n)\) time.
    Testing \(t\) independent bases thus has bit complexity \(O(t \log^3 n)\).
\end{remark}
